% source: https://tex.stackexchange.com/a/759208 (question 759198, score 12, block 1)
\documentclass[border=5pt]{standalone}% compile with lualatex only
\usepackage[svgnames]{xcolor}
\usepackage[3d]{luadraw}%https://github.com/pfradin/luadraw
\usepackage{fourier-otf}
%https://tex.stackexchange.com/questions/759198/improve-of-solid-efffect-of-stack-of-unit-cube-legos-solid-draw-with-hole-prot
\begin{document}
\begin{luadraw}{name=cube_lego}
local g = graph3d:new{viewdir=perspective("yz",0.5,135)}

local Dunitcube = function(O, color, prot)
    -- O = reference point
    -- prot = "top" or "left" or "front"
    local cyl = function(A,V,r)
        g:Dpoly(cylinder(A,V,r,45),{color=color, contrast=1.5,backcull=true, mode=mShadedOnly})
    end
    local r, h = 0.2, 0.1
    local circles_opt1 = 'draw=none,fill=black'
    local circles_opt2 = 'draw=none,fill='..color.."!80!black"
    g:Dpoly( parallelep(O,vecJ, -vecI,vecK), {color=color, mode=mShadedOnly, backcull=true})
    --front
    if prot == "front" then cyl(O+M(0,0.5,0.5), h*vecI, r)
    else g:Dcircle3d( O+M(0,0.54,0.5), r, vecI,circles_opt1); g:Dcircle3d(O+M(0,0.5,0.5),r,vecI,circles_opt2)
    end
    -- left
    if prot == "left" then cyl(O+M(-0.5,0,0.5), -h*vecJ, r)
    else g:Dcircle3d( O+M(-0.55,0,0.5), r, vecJ,circles_opt1); g:Dcircle3d(O+M(-0.5,0,0.5),r,vecJ,circles_opt2)
    end
    -- top
    if prot == "top" then cyl(O+M(-0.5,0.5,1), h*vecK, r)
    else g:Dcircle3d( O+M(-0.54,0.52,1), r, vecK,circles_opt1); g:Dcircle3d(O+M(-0.5,0.5,1),r,vecK,circles_opt2)
    end
end

-- some examples
Dunitcube(Origin,"SteelBlue","left"); Dunitcube(Origin-vecJ,"SeaGreen","left")
Dunitcube(Origin+vecK,"Gold",nil); Dunitcube(Origin-2*vecJ,"Beige","left"); 
Dunitcube(Origin-3*vecJ,"Crimson","left")
local A = 3*vecJ-3*vecK
Dunitcube(A,"SteelBlue","top"); Dunitcube(A+vecK,"SeaGreen","top")
Dunitcube(A+2*vecK,"Gold","top"); Dunitcube(A+3*vecK,"Beige","top")
Dunitcube(A+4*vecK,"Salmon","top"); Dunitcube(A+5*vecK,"cyan","top")
g:Dlabel3d("$5$", Origin-vecJ, {pos="S"}, "$6$", A+0.5*vecJ, {pos="S"})
g:Show()
\end{luadraw}
\end{document}
